\section{Related Work}
\subsection{Verifiable and Differentiable Physical Feedback}
RLVP post-trains language models to generate PDE solver code using continuous solution-accuracy and residual-consistency rewards~\cite{cai2026rlvp}. Its residual diagnostic scores the generated solution against a reference-based verifier. ISPO studies structured designs and the discrepancy between cheap and target evaluations; the proposed residual signal is evaluated on a cheaply obtained field for the same design.

PIRF fine-tunes diffusion models through differentiable physical rewards~\cite{yuan2025pirf}. GRADE propagates reward gradients through relaxed language-model generation~\cite{nel2025grade}. GRADE's shared-vocabulary reward interface differs from an external physical solver that consumes parsed numerical parameters. The numerical-head bridge considered here addresses this interface difference. Direct physical backpropagation and continuous physics rewards are established ingredients; their usefulness under heterogeneous evaluation cost is the subject of this study.

GI-PPO adaptively incorporates analytical gradients into policy optimization~\cite{son2023gippo}. HPO explicitly combines score-function and pathwise gradients for hybrid discrete--continuous action spaces~\cite{alvo2026hpo}. These works motivate matched gradient-only and mixed-gradient comparisons. ISPO adds a target evaluator whose feedback may be expensive or partially exposed, requiring a correction and allocation rule across evaluation interfaces.

\subsection{Multi-Fidelity and Gradient Control Variates}
MFPG constructs policy-gradient control variates from correlated low- and high-fidelity samples~\cite{liu2025mfpg}. The October 1, 2026 MFPG-PPO extension addresses correlation loss and variance inflation in modern actor--critic training and includes budget-aware allocation~\cite{liu2026mfpgppo}. It is the principal multi-fidelity baseline. MFPG-PPO explicitly studies the correlation--bias tradeoff caused by re-synchronizing low-fidelity rollouts, and its allocation analysis targets an idealized scalar-surrogate variance. The present information-selection study separately measures parameter-gradient variability and the cost of physical information interfaces. The present problem evaluates the same design at both fidelities, so coupling does not require maintaining matched physical trajectories. The proposed distinction is the use of objective-aware residual and sensitivity signals inside the auxiliary update and acquisition rule.

REBAR and RELAX combine score-function estimation with differentiable control variates for discrete or black-box objectives~\cite{tucker2017rebar,grathwohl2018relax}. Q-Prop combines an analytical critic-based term with a Monte Carlo correction~\cite{gu2017qprop}. These precedents establish the broader idea of using differentiable auxiliary information without optimizing the auxiliary objective alone. The conditional correction below is presented as a candidate construction for the physical design setting, rather than as a new general control-variate identity.


RLOO provides a sequence-level leave-one-out policy-gradient baseline~\cite{ahmadian2024rloo}. ABC revisits learned advantage/value control variates for RLVR~\cite{pan2026abc}, and ABC-Align already studies prediction-powered LLM post-training with adaptive bias control~\cite{frankel2026abcalign}. PPI++ and active statistical inference supply related prediction/correction and acquisition principles~\cite{angelopoulos2024ppipp,zrnic2024active}. These works motivate direct learned-feedback and nonphysical-uncertainty baselines. ISPO's proposed distinction is the physical information available from evaluator states, residuals and adjoints, combined with matched hard designs and complete information cost accounting.

\subsection{Residual-Based Objective Error Estimation}
Adjoint-weighted residual correction is a classical technique for estimating output functionals from approximate PDE states~\cite{pierce2000adjoint}. Wang and Hu already place dual-weighted-residual mesh adaptation and functional calculation inside an RL airfoil-design evaluator~\cite{wang2024mechanism}. Their solver workflow includes fine-mesh solution steps. The present question is whether partly purchased physical information can improve post-training updates and evaluation allocation without requiring a full target solve for every design.

NeuroForge studies residual-based trust for neural CFD surrogates and distinguishes error detection from state correction~\cite{jabbary2026neuroforge}. This distinction is important here: a residual is calibrated against target objective error before it controls acquisition. Raw residual magnitude is not assumed to be a universal certificate for physical quality. The prospective contribution is a coherent and physically grounded combination of target correction, physical error information, and cost-aware post-training, with each numerical and learning ingredient attributed to prior work.


\subsection{Language-Guided Design and Digital Twins}
ChatTwin automates data-center digital-twin scene descriptions through a conversational language interface~\cite{li2023chattwin}. Phythesis combines language-driven layout proposals with physics-guided evolutionary refinement for simulation-ready data-center scenes~\cite{li2025phythesis}. These works establish the design application and motivate a named inference-time comparison. ISPO asks how the proposal policy itself can be post-trained using a corrected allocation of heterogeneous physical feedback across tasks.

Scientific Generative Agent uses language-model proposals for discrete scientific structure and differentiable simulations for continuous optimization~\cite{ma2024sga}. This is a close precedent for the discrete--continuous division of work. The post-training question concerns parameter updates to a reusable proposal policy and their fidelity/cost accounting. The aerodynamic protocol uses an established adjoint framework such as DAFoam~\cite{he2020dafoam}; solver development and conversational simulation tooling are attributed to their own contributions.
